% ============================================================
%  chapters/ch03-spherepop-visual-computation.tex
% ============================================================

\chapter{Spherepop as Visual Computation}

\chapterepigraph{%
  The parenthesis is a lie about space.\\
  It claims to enclose what it only annotates.\\
  The sphere actually encloses.
}{Flyxion, \textit{Geometry, Cognition, and the Transparency of Computation}}

\sidenote{App.\,C gives\\the full formal\\treatment}

Standard mathematical notation is a historical accident. The bracket
$(1 + (3 \times 2^2))$ encodes a spatial relationship — nesting — using
a pair of punctuation marks that have no spatial extent. The notation is
purely symbolic; it \emph{represents} containment without \emph{being}
containment.

Spherepop takes the spatial interpretation seriously. Expressions are
regions. Containment is actual enclosure. Evaluation is collapse. The
result is a notation system in which the \emph{picture} of an expression
and the \emph{meaning} of an expression are the same thing.

% ── 3.1 ──────────────────────────────────────────────────────
\section{Origins of Spherepop}

\begin{remark}
  \textbf{The plenum interpretation of evaluation.}
  Standard notation treats an expression as a string of symbols in
  empty syntactic space. Spherepop treats it as a physical medium —
  a structured plenum of nested constraint regions. This is not
  merely aesthetic. In the relaxing-spring cosmology of Part~V,
  the universe is itself a plenum (the lamphrodyne field) that
  ``evaluates'' by reducing constraint violation: high-constraint
  early states collapse toward lower-constraint equilibria, leaving
  persistent structures as their residue. Spherepop evaluation is
  the discrete, finite-depth version of this same dynamics.
  The pop operator at the innermost sphere is the computational
  analogue of a quantum event (Chapter~6) — irreversible, local,
  and constraint-reducing.
\end{remark}


The Spherepop formalism arose from a dissatisfaction with the gap between
expression and visualization. In conventional notation:

\begin{itemize}
  \item The tree structure of an expression must be \emph{inferred}
        from flat text.
  \item Precedence rules are \emph{remembered}, not \emph{seen}.
  \item Evaluation order is \emph{computed mentally}, not perceived
        directly.
\end{itemize}

A child learning arithmetic must perform three separate cognitive operations:
parse the linear string into a tree, remember the precedence rules, and then
evaluate bottom-up. Each of these is non-trivial. The notation makes
computation harder than it needs to be.

\begin{historynote}
  Lisp (1958) was the first major programming language to make expression
  trees visible in syntax. The expression $1 + 3 \times 4$ becomes
  \texttt{(+ 1 (* 3 4))} in Lisp — a parenthesized tree that matches
  the evaluation structure exactly. Lisp programmers sometimes call
  their parentheses ``hugs'': each open paren hugs its contents.
  Spherepop replaces the linear hugging with actual spatial enclosure,
  taking Lisp's structural insight one step further into geometry.
  The relationship between Spherepop and Lisp is roughly the relationship
  between a map and a legend: Spherepop makes the spatial structure of
  computation directly visible rather than encoding it symbolically.
\end{historynote}

% ── 3.2 ──────────────────────────────────────────────────────
\section{Nested Bubbles}

A Spherepop expression is a planar diagram of nested closed regions —
\emph{bubbles}. Each bubble contains either:
\begin{enumerate}
  \item an atomic value (a number, variable, or constant), or
  \item an operator together with its argument bubbles.
\end{enumerate}

The spatial containment \emph{is} the grouping. No parentheses are needed
because the containment relation is directly visible.

\begin{definition}{Spherepop Diagram}{sp-diagram}
  A \emph{Spherepop diagram} is a planar embedding of a rooted ordered
  tree into $\mathbb{R}^2$ such that:
  \begin{itemize}
    \item each node $v$ of the tree corresponds to a closed disk
          $D_v \subset \mathbb{R}^2$,
    \item if $u$ is the parent of $v$ in the tree, then $D_v \subset
          D_u^\circ$ (strict interior),
    \item sibling disks are pairwise disjoint:
          $D_{v_1} \cap D_{v_2} = \emptyset$ for distinct siblings,
    \item the label of node $v$ is written inside $D_v$ but outside
          all children of $v$.
  \end{itemize}
\end{definition}

\begin{figure}[h]
  \centering
  \begin{tikzpicture}[scale=0.72]
    \tikzset{
  bubble/.style={circle, draw, thick, minimum size=7mm, inner sep=0},
  poparrow/.style={-{Stealth[length=7pt]}, RSVPaccent, very thick},
  steplab/.style={RSVPdeep, font=\small\bfseries\sffamily},
  nodelab/.style={RSVPdeep, font=\footnotesize},
  collapselab/.style={RSVPaccent!80, font=\footnotesize\itshape},
}
% ============================================================
%  diagrams/spherepop-eval-3d.tex
%  Spherepop evaluation sequence as nested 3D bubbles
%  Shows: 1 + 3×2² evaluated step by step
% ============================================================
  bubble/.style={circle, draw, thick, minimum size=7mm, inner sep=0},
  poparrow/.style={-{Stealth[length=7pt]}, RSVPaccent, very thick},
  steplab/.style={RSVPdeep, font=\small\bfseries\sffamily},
  nodelab/.style={RSVPdeep, font=\footnotesize},
  collapselab/.style={RSVPaccent!80, font=\footnotesize\itshape},
% ── STEP 0: Full nested expression 1 + 3×(2²) ──────────────
\begin{scope}[xshift=0cm, yshift=0cm]
  \node[steplab] at (0, 2.8) {Step 0: Initial expression};

  % Outermost bubble (the + operation)
  \draw[RSVPdeep, thick] (0,0) circle (2.4cm);
  \node[nodelab] at (-1.3, 1.5) {$+$};

  % Left atom: 1
  \draw[RSVPmid, thick] (-1.2, 0) circle (0.55cm);
  \node[nodelab] at (-1.2, 0) {$1$};

  % Middle bubble: × operation
  \draw[RSVPmid, thick] (0.8, 0) circle (1.4cm);
  \node[nodelab] at (-0.1, 0.8) {$\times$};

  % Left of ×: atom 3
  \draw[RSVPpurple, thick] (0.1, 0) circle (0.45cm);
  \node[nodelab] at (0.1, 0) {$3$};

  % Right of ×: inner bubble for 2²
  \draw[RSVPaccent, thick] (1.5, 0) circle (0.7cm);
  \node[nodelab] at (1.3, 0.35) {\footnotesize$\hat{\phantom{.}}$};
  % atom 2
  \draw[RSVPaccent!70, thick] (1.25, -0.1) circle (0.28cm);
  \node[font=\tiny] at (1.25,-0.1) {$2$};
  % atom 2 (exponent)
  \draw[RSVPaccent!70, thick] (1.7, 0.25) circle (0.22cm);
  \node[font=\tiny] at (1.7,0.25) {$2$};

  % Pop target indicator
  \draw[RSVPaccent, dashed, thick, opacity=0.6]
    (1.5, 0) circle (0.78cm);
  \node[collapselab] at (1.5, -1.05) {pop here};
\end{scope}

% ── Arrow 1 ─────────────────────────────────────────────────
\draw[poparrow] (2.8, 0) -- (4.2, 0)
  node[midway, above, collapselab] {$\bullet_1$};

% ── STEP 1: 2² evaluated → 4 ─────────────────────────────────
\begin{scope}[xshift=5.2cm, yshift=0cm]
  \node[steplab] at (0, 2.8) {Step 1: $2^2 \to 4$};

  \draw[RSVPdeep, thick] (0,0) circle (2.4cm);
  \node[nodelab] at (-1.3, 1.5) {$+$};

  \draw[RSVPmid, thick] (-1.2, 0) circle (0.55cm);
  \node[nodelab] at (-1.2, 0) {$1$};

  \draw[RSVPmid, thick] (0.8, 0) circle (1.4cm);
  \node[nodelab] at (-0.1, 0.8) {$\times$};

  \draw[RSVPpurple, thick] (0.1, 0) circle (0.45cm);
  \node[nodelab] at (0.1, 0) {$3$};

  % 2² has been evaluated: now a single node = 4
  \draw[RSVPaccent, very thick, fill=RSVPaccent!15] (1.5, 0) circle (0.5cm);
  \node[nodelab, font=\small\bfseries] at (1.5, 0) {$4$};

  % Pop next target: ×
  \draw[RSVPaccent, dashed, thick, opacity=0.6]
    (0.8, 0) circle (1.48cm);
  \node[collapselab] at (0.8, -1.55) {pop here};
\end{scope}

% ── Arrow 2 ─────────────────────────────────────────────────
\draw[poparrow] (8.0, 0) -- (9.4, 0)
  node[midway, above, collapselab] {$\bullet_2$};

% ── STEP 2: 3×4 → 12 ─────────────────────────────────────────
\begin{scope}[xshift=10.4cm, yshift=0cm]
  \node[steplab] at (0, 2.8) {Step 2: $3{\times}4 \to 12$};

  \draw[RSVPdeep, thick] (0,0) circle (2.4cm);
  \node[nodelab] at (-1.3, 1.5) {$+$};

  \draw[RSVPmid, thick] (-1.2, 0) circle (0.55cm);
  \node[nodelab] at (-1.2, 0) {$1$};

  % 3×4 evaluated → 12
  \draw[RSVPmid, very thick, fill=RSVPmid!15] (0.8, 0) circle (0.6cm);
  \node[nodelab, font=\small\bfseries] at (0.8, 0) {$12$};

  % Pop next: outer +
  \draw[RSVPaccent, dashed, thick, opacity=0.6]
    (0,0) circle (2.48cm);
  \node[collapselab] at (0, -2.55) {pop here};
\end{scope}

% ── Arrow 3 (below, wrapping) ────────────────────────────────
\draw[poparrow] (10.4, -3.1) -- (10.4, -4.0) -- (5.4, -4.0) -- (5.4, -3.2);
\node[collapselab] at (7.9, -4.3) {$\bullet_3$};

% ── STEP 3: 1+12 → 13 (final) ────────────────────────────────
\begin{scope}[xshift=5.4cm, yshift=-6.4cm]
  \node[steplab] at (0, 2.4) {Step 3: $1{+}12 \to 13$ \quad (collapsed)};

  % Final result: single atom
  \draw[RSVPdeep, very thick, fill=RSVPdeep!10] (0,0) circle (0.85cm);
  \node[RSVPdeep, font=\Large\bfseries] at (0,0) {$13$};

  \node[collapselab, font=\small] at (0, -1.2)
    {evaluation complete — depth 0};
\end{scope}

% ── Depth legend ─────────────────────────────────────────────
\begin{scope}[xshift=0cm, yshift=-5.5cm]
  \node[draw=RSVPdeep!30, fill=RSVPgray, rounded corners=3pt,
        inner sep=6pt, font=\scriptsize, anchor=north west] at (0,0) {
    \begin{tabular}{ll}
      \color{RSVPdeep}\rule{10pt}{1.5pt} & Outer sphere (depth 0–1)\\
      \color{RSVPmid}\rule{10pt}{1.5pt}  & Mid sphere (depth 1–2)\\
      \color{RSVPaccent}\rule{10pt}{1.5pt}& Inner sphere (depth 2–3)\\
      $\bullet_n$ & Pop at step $n$
    \end{tabular}
  };
\end{scope}

  \end{tikzpicture}
  \caption{Spherepop evaluation of $1 + 3 \times 2^2 = 13$.
    Each frame shows the expression after one pop. The innermost
    bubble collapses first, reducing depth at each step. The
    complexity functional $C(E)$ decreases $3 \to 2 \to 1 \to 0$.}
  \label{fig:sp-eval-ch3}
\end{figure}

% ── 3.3 ──────────────────────────────────────────────────────
\section{Reading Order}

A conventional mathematical expression is read left-to-right, with
precedence rules overriding the linear order. A Spherepop diagram is
read \emph{inside-out}: innermost bubbles are read first, outermost last.

This reading order is not arbitrary. It is the \emph{only} reading order
consistent with the semantics: you cannot evaluate an outer expression
without first knowing the values of its inner arguments.

\begin{proposition}{Unique Canonical Reading Order}{reading-order}
  In any Spherepop diagram, inside-out reading (depth-first, innermost
  first) is the unique reading order consistent with compositional
  semantics.
\end{proposition}

\begin{proof}
  This follows directly from Proposition~\ref{prop:depth-order} of
  Chapter~2: in any containment-based evaluation system, outer regions
  depend on the values of inner regions, so inner regions must be
  evaluated first.
\end{proof}

The remarkable consequence is that Spherepop notation contains its own
reading instructions. A child presented with a Spherepop diagram for
the first time, and told only that ``inner bubbles evaluate first,'' has
everything they need to compute any arithmetic expression correctly —
without memorizing a single precedence rule.

% ── 3.4 ──────────────────────────────────────────────────────
\section{The Pop Operator}

The fundamental operation of Spherepop evaluation is the \emph{pop}: the
collapse of an innermost bubble to a single value.

\begin{definition}{Pop}{pop-ch3}
  A bubble $B$ is \emph{poppable} if it contains no child bubbles —
  i.e., it is an innermost bubble containing only atoms and an operator.
  The \emph{pop} of $B$ replaces $B$ with the result of applying its
  operator to its atomic arguments.
\end{definition}

\begin{example}
  In the diagram for $1 + (3 \times (2 \uparrow 2))$:
  \begin{enumerate}
    \item First pop: the innermost bubble $(2 \uparrow 2)$ pops to $4$.
    \item Second pop: the bubble $(3 \times 4)$ pops to $12$.
    \item Third pop: the outermost bubble $(1 + 12)$ pops to $13$.
  \end{enumerate}
  Three pops, three depth reductions, final result $13$.
\end{example}

The pop operator has two important special cases:

\begin{definition}{Refuse and Bind}{refuse-bind}
  A bubble \emph{refuses} when it cannot be popped — either because
  it contains unbound variables, because its operator is undefined on
  its arguments, or because a protection marker prevents evaluation.
  A bubble is \emph{bound} when its free variables are assigned values,
  enabling a previously refusing bubble to pop.
\end{definition}

Refusal is not failure — it is \emph{suspended evaluation}. A refusing
bubble waits until the conditions for its pop are satisfied. This is the
Spherepop model of a closure, a lazy value, or a blocked computation.

% ── 3.5 ──────────────────────────────────────────────────────
\section{Collapse as Evaluation}

The full evaluation of a Spherepop diagram is a sequence of pops that
reduces the diagram to a single atomic bubble. This process is
\emph{collapse} in the topological sense: the nested structure contracts
level by level until no further nesting remains.

\begin{theorem}{Spherepop Completeness}{sp-completeness}
  Every finite Spherepop diagram with no refusing bubbles and no
  unbound variables collapses to a single atomic bubble in finitely
  many pops.
\end{theorem}

\begin{proof}
  By induction on the depth $d$ of the diagram. For $d = 0$, the
  diagram is already atomic. For $d \geq 1$, at least one innermost
  bubble exists (depth-$d$ bubble with no children). This bubble pops,
  reducing depth by at least one. By induction the resulting depth-$(d-1)$
  diagram eventually collapses.
\end{proof}

This theorem is why Spherepop evaluation terminates. The depth of a diagram
is a natural number; each pop reduces it; natural numbers have no infinite
descending chains.

% ── 3.6 ──────────────────────────────────────────────────────
\section{Visual Examples from Arithmetic}

The power of Spherepop notation becomes clear in examples where conventional
notation is genuinely confusing.

\begin{example}[title={Order of operations}]
  The expression $2 + 3 \times 4$ is ambiguous in natural language
  (``two plus three, times four''? or ``two plus, three times four''?)
  and requires precedence rules in conventional notation. In Spherepop,
  the grouping is visible: if $\times$ is a smaller inner bubble and
  $+$ is the outer bubble, then $3 \times 4 = 12$ pops first and
  $2 + 12 = 14$ follows. The diagram is unambiguous without any rules.
\end{example}

\begin{example}[title={Nested exponentiation}]
  Consider $2^{2^3}$. In conventional notation this is ambiguous:
  is it $(2^2)^3 = 64$ or $2^{(2^3)} = 256$? In Spherepop, the
  nesting is spatial: whichever exponentiation bubble contains the other
  evaluates last. The conventional convention (right-to-left for towers)
  corresponds to making the rightmost exponent the innermost bubble.
\end{example}

\begin{example}[title={Function composition}]
  $f(g(h(x)))$ in Spherepop is three nested bubbles: $h(x)$ inside
  $g(\cdot)$ inside $f(\cdot)$. The evaluation order — $h$ then $g$
  then $f$ — is directly visible. The notation makes it impossible to
  confuse $f \circ g \circ h$ with $h \circ g \circ f$.
\end{example}

% ── 3.7 ──────────────────────────────────────────────────────
\section{Relationship to Lisp Trees and Expression Parsing}

\sidenote{Lisp: McCarthy,\\1958; see also\\Church's $\lambda$-calculus,\\1932}

Spherepop diagrams are visual renderings of the \emph{abstract syntax trees}
(ASTs) that all programming language parsers construct internally. When a
compiler parses $1 + 3 \times 4$, the first thing it does is build the tree:

\medskip
\begin{center}
\begin{tikzpicture}[
  level distance=1.2cm,
  sibling distance=2.2cm,
  every node/.style={circle, draw=RSVPdeep, fill=RSVPlight,
    minimum size=7mm, font=\small}
]
\node {$+$}
  child { node {$1$} }
  child { node {$\times$}
    child { node {$3$} }
    child { node {$4$} }
  };
\end{tikzpicture}
\end{center}
\medskip

A Spherepop diagram is this same tree, but rendered spatially: each internal
node becomes a bubble enclosing its children. The tree structure is the same;
only the rendering changes.

This observation has a deep implication: \emph{every programming language
with well-defined operator precedence already implicitly uses Spherepop
evaluation}. The parser builds the tree; the evaluator collapses it
depth-first. Spherepop makes this hidden structure visible.

\begin{remark}
  The relationship between Spherepop and conventional notation is exactly
  the relationship between a subway map and a street map. The street map
  shows the geography; the subway map shows the structure. Neither is
  more correct. They serve different cognitive purposes. Conventional
  notation is optimized for writing linearly (a typewriter constraint
  inherited from print). Spherepop is optimized for visual understanding
  of structure.
\end{remark}

% ── 3.8 ──────────────────────────────────────────────────────
\section{Spherepop Beyond Arithmetic}

The Spherepop formalism extends naturally beyond arithmetic to any domain
organized by containment and evaluation.

\begin{itemize}
  \item \textbf{Logic:} propositions are bubbles; logical connectives
        are operators; truth evaluation is collapse. The Hasse diagram
        of Chapter~1 and Appendix~D is a Spherepop structure over the
        Boolean lattice.
  \item \textbf{Type theory:} types are constraint boundaries; terms
        inhabit the interior; type-checking is evaluating whether a
        term pops to an admissible value.
  \item \textbf{Grammar:} syntactic constituents are bubbles; grammatical
        rules are operators; parsing is the collapse of a word sequence
        to a phrase structure tree.
  \item \textbf{Biology:} cells are bubbles containing organelles;
        organisms are bubbles containing cells; ecosystems are bubbles
        containing organisms. Metabolism is evaluation; death is collapse
        to a terminal state.
  \item \textbf{Cosmology:} the universe in the RSVP framework is an
        accessibility field; structures (galaxies, stars, planets, minds)
        are nested pockets of concentrated density; evolution of the
        universe is the progressive collapse of high-entropy regions
        into stable low-entropy structures.
\end{itemize}

\begin{conceptaside}{The Universality of Nested Evaluation}
  The claim that Spherepop generalizes to all these domains is not a
  metaphor. It is a precise mathematical claim: all these systems are
  constraint spaces $(X, A, \kappa)$ with nesting sequences and collapse
  maps $\mathcal{K}$. The formal apparatus of Appendix~C applies to all
  of them.

  What varies between domains is not the mathematics but the
  \emph{interpretation}: what counts as a bubble, what counts as an
  operator, what counts as a pop, and what the admissible region $A$
  is. The structure is universal; the semantics are domain-specific.

  This is a general principle in mathematics: the same formal structure
  appears in radically different domains because those domains are all
  special cases of the same underlying abstraction. Group theory,
  topology, category theory — each was discovered in a specific context
  and turned out to be universal. Spherepop is, we argue, a similar
  discovery: the formal structure of nested evaluation, found first in
  arithmetic, is universal.
\end{conceptaside}

% ── Exercises ────────────────────────────────────────────────
\section*{Exercises}

\begin{exercise}
  Draw the Spherepop diagram for each of the following expressions.
  Then evaluate each by performing the pops in order, showing the
  diagram after each pop.
  \begin{enumerate}[label=(\alph*)]
    \item $(2 + 3) \times (4 - 1)$
    \item $2^{(3^2)}$
    \item $\sqrt{(3^2 + 4^2)}$
  \end{enumerate}
\end{exercise}

\begin{exercise}
  A Spherepop diagram with $n$ poppable bubbles requires exactly $n$
  pops to fully evaluate. Prove this. (Hint: show that each pop
  removes exactly one internal node from the AST.)
\end{exercise}

\begin{exercise}
  Formalize the claim that ``inner bubbles evaluate first'' is the
  unique reading convention consistent with compositional semantics.
  Specifically: define \emph{compositional semantics} for a Spherepop
  diagram, and then prove that any reading order other than inside-out
  produces undefined behavior for some diagram.
\end{exercise}

\begin{exercise}[title={Extension}]
  Extend the Spherepop formalism to handle \emph{infinite} nesting
  sequences — diagrams where some bubble contains infinitely many
  sub-bubbles. What additional conditions are needed to guarantee
  that evaluation terminates? Give an example of an infinite Spherepop
  diagram that evaluates and one that does not.
\end{exercise}
